\chapter{Tình Bạn \& Nghĩa Khí (Friendship \& Loyalty)}

\begin{center}
  \textit{\color{truyenblue}``A friend in need is a friend indeed.''}\\
  \textit{(Bạn trong lúc hoạn nạn mới là bạn đích thực.)}
\end{center}

\ngancach

\section{The Shared Umbrella (Chiếc Ô Chung)}
\begin{truyen}{Tình Bạn}{Sẻ Chia}
Two \textbf{friends} walk in the \textbf{rain}. One has no \textbf{umbrella}. The other \textbf{shares} his, and they both get a little \textbf{wet}.

\textit{Hai người bạn đi dưới mưa. Một người không có ô. Người kia chia sẻ chiếc ô của mình, và cả hai đều bị ướt một chút.}

\begin{baihoc}
Tình bạn đích thực là sẵn sàng chia sẻ khó khăn cùng nhau.
\end{baihoc}
\end{truyen}

\section{The Silent Support (Sự Ủng Hộ Lặng Thầm)}
\begin{truyen}{Tình Bạn}{Cảm Thông}
A boy is \textbf{sad} and sits \textbf{alone}. His friend sits \textbf{quietly} next to him. Words are not \textbf{needed}.

\textit{Một cậu bé buồn bã và ngồi một mình. Bác bạn cậu ngồi lặng lẽ bên cạnh. Không cần đến lời nói.}

\begin{baihoc}
Đôi khi, sự hiện diện lặng thầm còn quý giá hơn ngàn lời an ủi.
\end{baihoc}
\end{truyen}

\section{The Broken Pencil (Chiếc Bút Chì Gãy)}
\begin{truyen}{Tình Bạn}{Giúp Đỡ}
A student \textbf{breaks} his \textbf{pencil} before a \textbf{test}. His \textbf{classmate} snaps his own pencil in \textbf{half} to share.

\textit{Một học sinh làm gãy bút chì trước giờ kiểm tra. Bạn cùng lớp của cậu bẻ đôi chiếc bút chì của mình để chia sẻ.}

\begin{baihoc}
Một hành động nhỏ của sự sẻ chia có thể mang lại ý nghĩa lớn lao.
\end{baihoc}
\end{truyen}

\section{The True Winner (Ngôi Vị Đích Thực)}
\begin{truyen}{Tình Bạn}{Đồng Hành}
Two girls \textbf{run} in a \textbf{race}. One falls down. The other \textbf{stops}, helps her up, and they cross the \textbf{finish} line together.

\textit{Hai cô gái chạy trong một cuộc đua. Một người bị ngã. Người kia dừng lại, giúp cô ấy đứng lên, và họ cùng nhau vượt qua vạch đích.}

\begin{baihoc}
Chiến thắng đích thực không phải là về đích đầu tiên, mà là cùng nhau vượt qua thử thách.
\end{baihoc}
\end{truyen}

\section{The Lost Dog (Chú Chó Bị Lạc)}
\begin{truyen}{Tình Bạn}{Quan Tâm}
A boy \textbf{loses} his dog. All his \textbf{friends} spend the whole \textbf{day} searching the \textbf{neighborhood} until they find him.

\textit{Một cậu bé bị mất con chó của mình. Tất cả bạn bè của cậu dành cả ngày tìm kiếm khắp khu phố cho đến khi họ tìm thấy nó.}

\begin{baihoc}
Những người bạn thực sự sẽ luôn sát cánh bên bạn trong những lúc hoạn nạn.
\end{baihoc}
\end{truyen}

\section{The Secret Keeper (Người Giữ Bí Mật)}
\begin{truyen}{Tình Bạn}{Tin Tưởng}
A girl \textbf{tells} her friend a \textbf{secret}. Even when they are \textbf{mad} at each other, the secret is never \textbf{told}.

\textit{Một cô gái kể cho bạn bè một bí mật. Ngay cả khi họ giận nhau, bí mật đó không bao giờ được tiết lộ.}

\begin{baihoc}
Sự tin tưởng và tôn trọng là nền tảng của một tình bạn lâu bền.
\end{baihoc}
\end{truyen}

\section{The Heavy Books (Những Cuốn Sách Nặng)}
\begin{truyen}{Tình Bạn}{Sẵn Sàng}
A boy \textbf{drops} his heavy \textbf{books}. Everyone \textbf{laughs}, except one \textbf{friend} who helps him \textbf{pick} them up.

\textit{Một cậu bé làm rơi những cuốn sách nặng của mình. Mọi người đều cười, ngoại trừ một người bạn giúp cậu nhặt chúng lên.}

\begin{baihoc}
Một người bạn tốt là người giang tay giúp đỡ khi người khác quay lưng đi.
\end{baihoc}
\end{truyen}

\section{The Honest Friend (Người Bạn Trung Thực)}
\begin{truyen}{Tình Bạn}{Chân Thành}
A boy \textbf{sings} badly on \textbf{stage}. His friend \textbf{tells} him the \textbf{truth} gently, so he can \textbf{practice} more.

\textit{Một cậu bé hát dở trên sân khấu. Bạn của cậu nhẹ nhàng nói sự thật, để cậu có thể luyện tập thêm.}

\begin{baihoc}
Lời nói thật lòng dù khó nghe cũng tốt hơn những lời nịnh bợ giả dối.
\end{baihoc}
\end{truyen}

\section{The Long Distance (Khoảng Cách)}
\begin{truyen}{Tình Bạn}{Gắn Kết}
Two friends move to \textbf{different} cities. They \textbf{write} letters \textbf{every} month to \textbf{keep} their friendship strong.

\textit{Hai người bạn chuyển đến các thành phố khác nhau. Họ viết thư mỗi tháng để giữ cho tình bạn của họ bền chặt.}

\begin{baihoc}
Khoảng cách địa lý không thể chia cắt được những tâm hồn thực sự gắn kết.
\end{baihoc}
\end{truyen}

\section{The Last Slice (Lát Bánh Cuối Cùng)}
\begin{truyen}{Tình Bạn}{Nhường Nhịn}
There is only one \textbf{slice} of pizza \textbf{left}. Two hungry \textbf{friends} insist the \textbf{other} should \textbf{eat} it.

\textit{Chỉ còn lại một lát bánh pizza. Hai người bạn đang đói cứ khăng khăng người kia nên ăn nó.}

\begin{baihoc}
Sự nhường nhịn và quan tâm là biểu hiện đẹp nhất của tình bạn chân thành.
\end{baihoc}
\end{truyen}
